\scanpage{72}
\uline{Th.} Les théories algébriques pour $\Delta$ sont celles qui sont 2-représentées par des $R\in\mathrm{Ob}\,\Delta$ admettant des \uline{petites familles \inserted{str.} cogénératrices}. Pour que $I\xrightarrow{\rho_0}R$ \deleted{\unsure{avec petite image dans|ayant petite image dans}} soit \inserted{$\Delta$-}basique, il f. et s. que son image soit \inserted{str. $\Delta$-}cogénératrice (c'est plus faible que \inserted{str.} cogénératrice !).
% slanted
\margin{passe bien dans la construction transfinie \unsure{si on prouve|d'en prouver} stabilité par $R[H]_\Delta$}

La nécessité dans la première assertion est claire par ce qui précède. Pour la suffisance, \deleted{on peut supp.} elle résulte de la deuxième assertion. Or soit $I'$ la catégorie \inserted{mun. d'un foncteur \uline{pl. fidèle} $I'\xrightarrow{\rho}R$,} telle que $\mathrm{Ob}\,I'=I$, $\mathrm{Ob}\,\rho=\varphi_0$. \unsure{Considérons|Considérons. On|\deleted{Considérons} On} \deleted{alors} \unsure{Sait|sait} que $T_{\widetilde I'}$ est une théorie algébrique de base $I$. Or $\widehat{I'^\circ}\longrightarrow R^\circ$ prolongeant $\rho:I'\longrightarrow R$ est ess. surjectif \inserted{et surjectif sur les Hom,} donc $\widetilde I'\longrightarrow R$ aussi. Donc il \deleted{suffit de} reste à rendre \deleted{\unsure{égales}} des \deleted{\unsure{iso}} flèches, donc rendre certaines flèches des iso.

\uline{Remarques} \textcircled{1} On \deleted{précise} \inserted{obtient} ici \deleted{à} la construction de tte théorie algébrique en les étapes suivantes : 1°) \unsure{choix|élts|objets} de base $I$ 2°) rajouter des flèches 3°) axiomes d'égalité sur des \inserted{composés de} flèches (on a maintenant $T_{\widetilde I'}$), 4°) \deleted{axiomes d'égalité sur d'autres \inserted{composés de} flèches, ou encore} \inserted{ou encore rendre inversibles des flèches}

On peut d'ailleurs condenser 3° et 4°) en une seule étape, \deleted{\unsure{3}} (imposer des axiomes successivement \unsure{revient à|rev. à} le faire tjrs en une seule fois, grâce au calcul de fractions) donc on donne $R$ \deleted{par} comme catégor. de fractions de $\widetilde I'=\widehat{I'^\circ}{}^\circ$ (d'un opposé de topos) par un ens. de flèches satisfaisant les conditions de saturation déjà dites.

\textcircled{2} Soit $R\in\mathrm{Ob}\,\Delta$. Les catégories quotients $R'\in\mathrm{Ob}\,\Delta$ déduites de $R$ \deleted{par} rendant inversibles des flèches \deleted{sont exactement celles qui} \inserted{correspondent 1--1 aux} ens. $\Sigma$ de flèches \deleted{$\Sigma$} telles que

a) $\Sigma$ stable par composition, par changement de base, contient les iso

b) $\Sigma$ stable par produits quelconques.

\textcircled{3} Supposons que $R$ ait une petite famille cogénératrice $I$. Alors \deleted{les} $\underline{\mathrm{Hom}}_{\Delta}(R,\mathrm{Ens})\simeq R^\circ$ [les foncteurs de $R$
% archive sheet number: 71
